\section{Prüfräder vor dem Versuch}
\label{sec:pruefraeder_vorher}

Die Prüfräder wurden vor den Laufversuchen nach dem Verfahren aus \secref{sec:dokumentation_pruefraeder} vermessen. Von den zehn Kernrädern der Laufversuche und den 26 Rädern der Fertigungsstreuung liegen 58 Dateien der Verzahnungsmessung vor, von den sieben Ritzeln der Prüfstandsserie zehn Dateien. In mehreren Runden wurden 21 Räder und drei Ritzel gemessen. Dazu kommen 47 Konturscans von 18 Rädern, je drei Scans des Prüfstandsritzels 95912 und des älteren Ritzels 86481 sowie Rauheitsspuren von 35 Rädern. Sechs Kernräder sind nicht gescannt, für die Ritzel fehlt die Rauheitsmessung.

\tabref{tab:kernraeder_gina} stellt die Kenngrößen der Kernräder denen der Fertigungsstreuung gegenüber. Je Rad zählt die letzte Messung. Die Toleranzen des Messprogramms sind die Werte der Verzahnungsqualität 7 nach DIN 3962 Teil 1 für Teilung, Rundlauf und Profil-Winkelabweichung \cite[S.~3]{DIN3962-1} und nach Teil 2 für die Flankenlinien-Winkelabweichung \cite[S.~1]{DIN3962-2}. Die Dateien nennen die Norm nicht. Sie führen aber deren Größen Teilungssprung und Teilungs-Spannenabweichung über ein Achtel des Umfangs, und die hinterlegten Werte stimmen mit den Tabellen der Norm überein. DIN 3962 ist durch DIN ISO 1328-1 ersetzt, deren Klassen anders gestuft sind \cite[Nationales Vorwort, S.~4 bis 5]{DINISO1328-1}. Die Zeichnung des Rades fordert die Klasse 10 nach DIN ISO 1328-1 (\secref{sec:pruefverzahnung}). Für $d = 52$\,mm, $m_{\mathrm{n}} = 1$\,mm und $b = 15$\,mm ergeben die Gleichungen (5) bis (12) der Norm mit ihren Rundungsregeln die Toleranzen $f_{\mathrm{pT}} = 31$\,µm, $F_{\mathrm{pT}} = 95$\,µm, $f_{\mathrm{H}\alpha\mathrm{T}} = 25$\,µm und $f_{\mathrm{H}\beta\mathrm{T}} = 32$\,µm \cite[S.~36 bis 37]{DINISO1328-1}. Der informative Anhang E nennt die Rundlauftoleranz $F_{\mathrm{rT}} = 85$\,µm \cite[S.~53]{DINISO1328-1}, der informative Anhang G die Toleranz des Teilungssprungs $f_{\mathrm{uT}} = 44$\,µm \cite[S.~60]{DINISO1328-1}. Gezählt werden die Überschreitungen beider Vorgaben.

Die Profil-Winkelabweichung $f_{\mathrm{H}\alpha}$ ist auf beiden Flanken positiv, ihr Mittel je Rad liegt zwischen $5{,}1$ und $18{,}0$\,µm. In jeder der 58 Dateien überschreiten einzelne Zähne die $7$\,µm des Messprogramms. Die Flankenlinien-Winkelabweichung $f_{\mathrm{H}\beta}$ wechselt von Zahn zu Zahn das Vorzeichen und erreicht im Betrag bis zu $107$\,µm. Ihr Mittel über die Zähne bleibt mit $5{,}5$\,µm links und $-4{,}6$\,µm rechts klein. Die Teilungs-Gesamtabweichung $F_{\mathrm{p}}$ überschreitet die $32$\,µm in 40 der 58 Dateien, der Rundlauf $F_{\mathrm{r}}$ die $22$\,µm in 55 Dateien. Das Rad 97345 erreicht mit $F_{\mathrm{p}}$ von $136{,}0$\,µm links und $122{,}9$\,µm rechts, $f_{\mathrm{p}}$ bis $43{,}1$\,µm und $F_{\mathrm{r}}$ von $56{,}9$\,µm die größten Werte der Kerngruppe. Das Rad 97361 folgt mit $80{,}7$ und $79{,}2$\,µm. In der Fertigungsstreuung liegen die vier Räder 97342, 97344, 97348 und 97369 mit $F_{\mathrm{p}}$ über $60$\,µm. Gegen die Klasse 10 der Zeichnung halten neun der zehn Kernräder Teilung und Rundlauf ein, nur das Rad 97345 überschreitet $F_{\mathrm{pT}}$ und $f_{\mathrm{pT}}$. In der Fertigungsstreuung liegen die Räder 97342, 97348 und 97369 über $F_{\mathrm{pT}}$, das Rad 97342 mit $88{,}7$\,µm auch über $F_{\mathrm{rT}}$. Die Profil- und die Flankenlinien-Gesamtabweichung bleiben an allen Rädern unter den $40$ und $48$\,µm der Klasse. Die Flankenlinien-Winkelabweichung einzelner Zähne, auf die Zahnbreite bezogen, überschreitet die $32$\,µm dagegen an acht der zehn Kernräder und an allen 26 Rädern der Fertigungsstreuung. Die Profil-Winkelabweichung einzelner Zähne liegt an einem Kernrad und drei Rädern der Streuung über $25$\,µm. Die Zeichnungsqualität ist damit in Teilung, Rundlauf und Gesamtabweichungen fast durchweg erreicht, in der Flankenlinien-Winkelabweichung einzelner Zähne nicht.

\begin{table}[htbp]
	\caption{Verzahnungsmessung der Räder vor dem Versuch: Mittel, Standardabweichung $s$ und Spanne über die Räder je Gruppe, Flankenwerte links / rechts, je Rad die letzte Messung. Toleranzen der Qualität 7 nach DIN 3962 \cite[S.~3]{DIN3962-1} \cite[S.~1]{DIN3962-2}, wie im Messprogramm hinterlegt (Spalte Q7), und der Klasse 10 nach DIN ISO 1328-1 \cite[S.~36 bis 37, 53 und 60]{DINISO1328-1}, wie die Zeichnung sie fordert (Spalte Q10).}
	\label{tab:kernraeder_gina}
	\centering
	\footnotesize
	\setlength{\tabcolsep}{3pt}
	\begin{tabular}{lllllllll}
		                       &         &          &            & \multicolumn{3}{l}{Kernräder, $n = 10$} & \multicolumn{2}{l}{Fertigungsstreuung, $n = 26$}                                                                         \\
		Größe                  & Einheit & Q7 & Q10  & Mittel                                  & $s$                                              & Spanne                    & Mittel              & $s$                 \\
		\hline
		\arrayrulecolor{black!30}
		$W_{\mathrm{k}}$       & mm      &          &            & $16{,}973$                              & $0{,}006$                                        & $16{,}962$ bis $16{,}984$ & $16{,}982$          & $0{,}004$           \\
		\hline
		$d_{\mathrm{a}}$       & mm      &          &            & $54{,}049$                              & $0{,}015$                                        & $54{,}014$ bis $54{,}068$ & $54{,}070$          & $0{,}013$           \\
		\hline
		$f_{\mathrm{p}}$       & µm      & $10$     & $31$ & $13{,}7$ / $15{,}6$                     & $8{,}9$ / $11{,}5$                               & $7{,}4$ bis $43{,}1$      & $14{,}9$ / $15{,}6$ & $10{,}3$ / $9{,}8$  \\
		\hline
		$f_{\mathrm{u}}$       & µm      & $12$     & $44$ & $15{,}3$ / $16{,}4$                     & $5{,}8$ / $5{,}7$                                & $9{,}5$ bis $31{,}1$      & $17{,}3$ / $18{,}8$ & $7{,}7$ / $10{,}5$  \\
		\hline
		$F_{\mathrm{p}}$       & µm      & $32$     & $95$ & $47{,}9$ / $48{,}0$                     & $36{,}8$ / $31{,}9$                              & $17{,}0$ bis $136{,}0$    & $45{,}6$ / $47{,}4$ & $28{,}4$ / $30{,}2$ \\
		\hline
		$F_{\mathrm{r}}$       & µm      & $22$     & $85$ & $37{,}2$                                & $11{,}6$                                         & $20{,}0$ bis $56{,}9$     & $43{,}3$            & $15{,}7$            \\
		\hline
		$f_{\mathrm{H}\alpha}$ & µm      & $7$      & $25$ & $10{,}3$ / $10{,}2$                     & $2{,}5$ / $1{,}6$                                & $5{,}1$ bis $13{,}3$      & $12{,}9$ / $12{,}7$ & $1{,}8$ / $2{,}2$   \\
		\hline
		$F_{\alpha}$           & µm      &          & $40$ & $18{,}1$ / $17{,}4$                     & $1{,}8$ / $1{,}8$                                & $15{,}4$ bis $21{,}6$     & $17{,}9$ / $17{,}3$ & $1{,}6$ / $1{,}3$   \\
		\hline
		$f_{\mathrm{H}\beta}$  & µm      & $11$     & $32$ & $5{,}5$ / $-4{,}6$                      & $6{,}4$ / $3{,}9$                                & $-12{,}5$ bis $13{,}7$    & $-5{,}1$ / $-1{,}8$ & $8{,}0$ / $8{,}1$   \\
		\hline
		$F_{\beta}$            & µm      &          & $48$ & $21{,}0$ / $20{,}3$                     & $6{,}4$ / $5{,}7$                                & $14{,}1$ bis $34{,}0$     & $21{,}4$ / $24{,}8$ & $5{,}0$ / $6{,}6$   \\
		\hline
		\arrayrulecolor{black}
	\end{tabular}
\end{table}

Die Wiederholmessungen an 21 Rädern grenzen die Messstreuung ab. Zwischen den Runden desselben Rades unterscheiden sich Zahnweite und Kopfkreisdurchmesser im Mittel um $0{,}005$\,mm und $0{,}010$\,mm, die Form- und Profilabweichungen um rund $2$\,µm. Bei $f_{\mathrm{H}\beta}$ links beträgt der mittlere Unterschied $10{,}4$\,µm und bei $F_{\mathrm{p}}$ $9{,}6$\,µm, die größten Spannweiten erreichen $33{,}5$\,µm und $32{,}3$\,µm. Aufbau und Ablauf waren in allen Runden gleich, je Runde wechselte der Bediener. Die Fertigungsstreuung der Steigungs- und Teilungsabweichungen ist deshalb zusammen mit dieser Messstreuung zu lesen. Die Teilungs-Sektorabweichung $F_{\mathrm{pk}}$ nach DIN ISO 1328-1 Anhang D.2 wurde aus den Teilungstabellen nachgerechnet, weil das Messprogramm unter Fpz/8 eine andere Spanne druckt. $F_{\mathrm{pk}}$ liegt auf 43 von 72 Flanken über $22$\,µm, die Programmgröße auf 41. Der Unterschied beider Größen erreicht $22$\,µm.

Die Konturscans der Räder bestätigen die Zahnform der Erzeugung und vermessen die Kopfkante. Der freie Evolventenfit mit dem Nenn-Grundkreis hatte die Profil-Winkelabweichung der Räder als radialen Versatz der Achse von im Median $52$\,µm zum Zahn hin aufgenommen. Dadurch fielen die Radien der Scans rund $50$\,µm zu klein und die Zahndicke rund $28$\,µm zu dünn aus. Mit der an den GINA-Kopfkreis desselben Teils gebundenen Achse beträgt die mittlere quadratische Abweichung (RMS) der Flanken $1{,}8$ bis $3{,}6$\,µm gegenüber $0{,}6$ bis $2{,}1$\,µm des freien Fits. Dieser Wert ist die reale Profilabweichung. Der Fußkreisdurchmesser der Scans weicht je Teil um $-31$ bis $+20$\,µm vom GINA-Wert ab, die Grundkreis-Zahndicke $s_{\mathrm{b}}$ liegt $5$ bis $15$\,µm über dem erzeugten Wert. Die Zahnweite $W_{\mathrm{k}}$ der Kernräder trifft mit $16{,}973$\,mm im Mittel den Wert der STplus-Erzeugung von $17{,}180$\,mm abzüglich $0{,}207$\,mm Zahndickenabmaß. Die Räder der Fertigungsstreuung liegen mit $16{,}982$\,mm $9$\,µm darüber. An der Kopfkante passt ein tangentialer Bogen in 162 von 163 Ecken besser als eine Fase. Sein Radius beträgt im Mittel je Teil $0{,}205$ bis $0{,}248$\,mm. Die Zeichnung des Rades verlangt eine Rundung r $0{,}2$, die STplus-Erzeugung trägt dagegen eine Fase von $45$°.

Die sieben Ritzel der Prüfstandsserie liegen in Teilung und Flankenform eng beieinander. $F_{\mathrm{p}}$ beträgt im Mittel $11{,}8$\,µm links und $10{,}0$\,µm rechts, $F_{\mathrm{r}}$ $13{,}6$\,µm und $f_{\mathrm{p}}$ $1{,}7$\,µm. Über die Toleranzen der Qualität 6 nach DIN 3962, die das Messprogramm hinterlegt, kommen nur wenige Werte. Die Zeichnung von 2022 fordert eine Qualität besser als 6 (\secref{sec:pruefverzahnung}). Der Rundlauf der Ritzel 95911 und 95915 liegt mit $16{,}2$ und $16{,}8$\,µm über $16$\,µm, $f_{\mathrm{H}\beta}$ einzelner Zähne des Ritzels 95920 erreicht $13{,}8$\,µm gegen $8$\,µm. Dazu kommt $f_{\mathrm{H}\alpha}$ der ersten Messrunde mit $5{,}2$ bis $7{,}0$\,µm, das in der zweiten Runde mit $-1{,}1$ bis $1{,}7$\,µm gemessen wurde. Die Kopfrücknahme liest das Messprogramm je Zahn mit $-22{,}2$ bis $-25{,}8$\,µm zwischen den Durchmessern $51{,}946$ und $52{,}880$\,mm, im Mittel $-23{,}3$\,µm links und $-23{,}7$\,µm rechts. Die Zahnweite streut mit einer Standardabweichung von $0{,}048$\,mm um $16{,}808$\,mm und liegt $4$\,µm unter dem erzeugten Wert von $16{,}812$\,mm. Zwischen den Runden schwankt sie beim Ritzel 95918 um $0{,}127$\,mm, Teilung und Rundlauf wiederholen innerhalb von $4{,}5$\,µm. Die Zeichnung von 2022 schreibt die Kopfrücknahme mit $25$\,µm linear ab dem Durchmesser $51{,}946$\,mm und die Zahnweite $16{,}822$\,mm mit dem Abmaß $-0{,}021$\,mm vor (\secref{sec:pruefverzahnung}).

Der Konturscan des Prüfstandsritzels 95912 ist die Gegenprobe des Achsverfahrens am Stahl. Die Bindung an den GINA-Kopfkreis verschiebt die Achse hier um $4$ bis $15$\,µm gegen $30$ bis $69$\,µm bei den gespritzten Rädern. Der Fußkreis stimmt mit $48{,}308$\,mm gegen $48{,}299$\,mm auf $9$\,µm. Die Zahndicke liegt $23$\,µm unter dem erzeugten Wert, wie es die Zahnweite $16{,}786$\,mm dieses Ritzels erwarten lässt. Die Rücknahme am Kopf liest der Scan mit $25{,}3$ bis $27{,}1$\,µm gegen $-22{,}7$ und $-23{,}4$\,µm der GINA-Messung. Die Kopfkante des Ritzels ist in \secref{sec:kopfkante} beschrieben. Das ältere Ritzel 86481 trägt eine Rücknahme von $20{,}8$\,µm und dient wegen seines um $0{,}17$\,mm tieferen Fußkreises nicht als Bezug (\secref{sec:pruefverzahnung}).

% TK erledigt (Kante nur noch in 4.2, 86481 nur noch in 3.1 beschrieben): Der Inhalt des Abschnitts ist schon mehr oder weniger an anderer Stelle genauso beschrieben. Bitte nur an einer Stelle das Ritzel in dem Ausmaß beschreiben. 

Die Rauheit der Zahnflanken wurde mit drei Spuren je Flanke gemessen, der Wert einer Flanke ist das Mittel der Spuren. Die Kernräder haben $R_{\mathrm{a}}$ von $0{,}81$\,µm links und $0{,}93$\,µm rechts bei Standardabweichungen von $0{,}32$ und $0{,}44$\,µm sowie $R_{\mathrm{z}}$ von $4{,}9$ und $5{,}8$\,µm. Die Räder der Fertigungsstreuung liegen mit $R_{\mathrm{a}}$ von $1{,}20$\,µm und $R_{\mathrm{z}}$ von $7{,}3$\,µm auf beiden Flanken darüber. Innerhalb der Fertigungsstreuung trennen sich zwei Gruppen nach Messtag und Bediener. Die am 12. und 21. Mai gemessenen Räder 97341 bis 97356 dieser Gruppe sowie 97371 liegen bei $R_{\mathrm{a}}$ von $1{,}2$ bis $1{,}9$\,µm. Die am 27. Mai unter dem Chargenlabel der Prüfverzahnung kst-E gemessenen Räder 97357 bis 97374 liegen bei $0{,}4$ bis $1{,}6$\,µm. Beide Gruppen stammen aus derselben Charge und demselben Messverfahren und wurden von zwei Bedienern aufgenommen. Ein anderer Unterschied besteht nicht. Je ein Protokoll der Verzahnungs- und der Rauheitsmessung steht als Beispiel in \anhref{sec:anhang_g}.

Für das FE-Modell mit gemessener Geometrie werden die Zahndicken aus den GINA-Zahnweiten übernommen, nicht aus den Scans. Die Grundlinie behält die in STplus erzeugte Zahndicke. Die Kopfrundung des Rades von $0{,}21$\,mm und die Kopfrücknahme des Ritzels gehen als örtliche Anpassungen ein, die Kopfkante des Ritzels bleibt nach \secref{sec:kopfkante} ein Parameter des Modells. Die Scans der gelaufenen Räder werden in \secref{sec:schadensbilder} mit dem Band aller Scans vor dem Versuch verglichen. Der Vergleich desselben Zahns setzt dieselbe Messebene voraus, deren Lage zur Stirnfläche des $15$\,mm breiten Zahns in den Dateien nicht festgehalten ist (beim Rad 97340 stimmt sie mit $85{,}95$\,mm überein). Die 30 Scans vom 23. Juni mit rund 400 Punkten bleiben aus der Kantenstatistik ausgeschlossen.

% TK erledigt (Absatz neu sortiert: Zahndicke, Kopfrundung und Rücknahme, Nachher-Scans, Kantenstatistik; Schliffbild und CT nur im Ausblick): Der letzte Abschnitt ist etwsa wirr formuliert und nimmt viele unterschiedliche Themen auf. Bitte klarer sortieren und aktualisieren auf den jetzigen Stand. Wenn da Fragen sind, bitte fragen.

% Herkunft:
% Absatz 1 (Umfang): architecture_decisions.md ADR-117 Consequences (10 Kernräder, 26 Räder der
%   Fertigungsstreuung, 7 Ritzel, 86481; Lücken) und Amendment 2026-10-12 Punkt 4 (Scan 95912);
%   80_output/messungen/gina/wheel/toleranzen_protokoll.md (58 Dateien) und gina/pinion/toleranzen_protokoll.md
%   (10 Dateien); gina/wheel/wiederholbarkeit.md (21 Teile), gina/pinion/wiederholbarkeit.md (3 Teile);
%   konturscan/wheel/bericht.md Tabelle "Je Scan" (47 Zeilen, 18 Teile) und konturscan/pinion/bericht.md
%   (je drei Scans 95912 und 86481); rauheit/wheel/gruppen.md (n 10 + 25); known_limits.md MEAS-04.
% Absatz 2 und Tabelle tab:kernraeder_gina: gina/wheel/tabelle_streuung.md (Mittel, s, min, max je Gruppe);
%   gina/wheel/toleranzen_protokoll.md Tabelle 1 (Toleranzen, Stufe) und Abschnitt "Überschreitungen"
%   (Werte und Dateien je Label; Einzelzeilen für 97345, 97361 -neu1, 97342, 97344, 97348, 97369 und
%   97344 rechts Zahn 1 = -107,2); known_limits.md MEAS-05 (2026-10-25 geschlossen: die Werte des Messprogramms
%   sind die Tabellen der DIN 3962 Teil 1 S. 2 und 3 und Teil 2 S. 1, an der Norm geprüft); Klasse 10 nach
%   DIN ISO 1328-1 aus gearcore.measurement.tolerances (Gl. (5) bis (12) S. 36 bis 37, Anhang E.1 S. 53, Anhang G.2
%   S. 60), Überschreitungen aus gina/wheel/toleranzen_protokoll.md Abschnitt Zeichnungsqualität (Kernräder: fp/Fp/fHa
%   nur 97345, fHb 8 von 10; Streuung: Fp 97342/97348/97369, Fr 97342 88,7, fHa 97342/97348/97364, fHb 26 von 26);
%   F_alphaT 40 und F_betaT 48 der Klasse 10 gegen die Spannen 15,4 bis 21,6 und 14,1 bis 34,0 der Tabelle.
% Absatz 3: gina/wheel/wiederholbarkeit.md ("Spannweite je Größe"); ADR-117 Amendment
%   2026-10-10 Punkt 4 und known_limits.md MEAS-08 (gleicher Aufbau, Bedienerwechsel);
%   gina/wheel/teilungssektoren.md (Kopf, Schlusszeile 41/43 von 72, 97348 rechts 78,1 gegen 100,1);
%   gina/wheel/fenster_zaehne.md (Spalte Rang: 97363 links, 97368 links Rang 1; Kernräder).
% Absatz 4: konturscan/wheel/bericht.md Kopf (Achsbindung, 52 µm, 50 µm, 28 µm), "Je Scan" (RMS Flanken
%   1,8 bis 3,6, frei 0,6 bis 2,1), "Je Teil" (d_f gegen GINA, s_b − erzeugt 5 bis 15, ρ Mittel 0,205 bis
%   0,248); gina/wheel/tabelle_streuung.md (W_k je Gruppe); ADR-117 Amendment
%   2026-10-12 Punkt 1 (17,180 − 0,207, +9 µm) und 2026-10-11 Punkt 2 (Zeichnung r 0,2, STplus 45°);
%   known_limits.md MEAS-11 (162 von 163 Ecken).
% Absatz 5: gina/pinion/tabelle_streuung.md (Mittel, s, W_k); gina/pinion/toleranzen_protokoll.md
%   (Toleranzen Stufe 6, Überschreitungen Fr, fHb 95920, fHa erste Runde); gina/pinion/wiederholbarkeit.md
%   (fHa zweite Runde, WK 95918 0,127, Fp/Fr max 4,4/4,5); gina/pinion/kopfruecknahme_ritzel.md (fKo je Zahn,
%   Beginn/Ende); ADR-117 Amendment 2026-10-12 Punkt 3 (16,808 gegen 16,812).
% Absatz 6: konturscan/pinion/bericht.md ("Je Scan" Versatz 15/4/5, ρ 0,004/0,006, Gerade 25,3 bis 27,1;
%   "Je Teil" d_f, s_b − erzeugt, 86481); ADR-117 Amendment 2026-10-12 Punkt 4 (Zahnweite 16,786,
%   Mikroskopbilder, 86481 anderer Schnitt, nicht Bezug von V2); known_limits.md MEAS-06.
% Absatz 7: rauheit/wheel/gruppen.md (Mittel, s, min, max); rauheit/wheel/tabelle.md
%   (Messtage und Chargenlabel je Rad); known_limits.md MEAS-04 (Ritzel ohne
%   Rauheit); ADR-117 Amendment 2026-10-11 Punkt 5 (Protokolle im Anhang).
% Absatz 8: ADR-117 Amendment 2026-10-12 Punkt 3 (Zahndicke aus Zahnweiten, V0/V1 STplus, Rundung und
%   Rücknahme örtlich); known_limits.md MEAS-06 (Kante bleibt Modellparameter), MEAS-10 (Messebene),
%   MEAS-03 (30 Scans vom 23.06.: 97363, 97364, 97367 bis 97374).
% Bildwunsch: 80_output/messungen/gina/wheel/zaehne_97345-neu.svg, Teilung und Rundlauf je Zahn des Kernrades
%   97345 mit dem Fenster der Zähne 36 bis 40 (größtes F_p der Kerngruppe); als Gegenstück zaehne_97366-neu.svg
%   (kleinstes F_p).
% Bildwunsch: GINA-Streuung der Räder (F_p, F_r, f_Halpha, W_k je Rad, Kerngruppe gegen Fertigungsstreuung),
%   noch zu zeichnen aus gina/wheel/tabelle_streuung.csv und den .mew-Werten (M6).
% Bildwunsch: Teilungssektoren Fpz/8 des Messprogramms gegen F_pk je Flanke aus gina/wheel/teilungssektoren.md,
%   noch zu zeichnen (M6).
% Bildwunsch: 80_output/messungen/konturscan/wheel/kopfkante_97340_97340-Z1.svg, Kopfrundung des Rades
%   (Messpunkte, scharfe Ecke, Bogen- und Fasenfit, Kernrad 97340, Zahn 1).
% Bildwunsch: 80_output/messungen/konturscan/wheel/ueberlagerung_97340_97340-Z1.svg, Zahnform des Rades gegen
%   die Erzeugung mit überhöhter Abweichung.
% Bildwunsch: 80_output/messungen/konturscan/pinion/ueberlagerung_pinion_rig_series.svg und
%   ueberlagerung_pinion_older.svg, Ritzel 95912 und 86481 gegen die Erzeugung (Kopfrücknahme, Fußkreis).
% Bildwunsch: Rauheitsgruppen R_a und R_z je Rad nach Gruppe, Messtag und Chargenlabel aus
%   rauheit/wheel/tabelle.md, noch zu zeichnen (ADR-117: bisher kein Rauheitsbild vorgesehen, Protokolle im
%   Anhang; Nutzer entscheidet).
